\chapter{Turning scripts into shared tools}\label{ch:tools}

Three colleagues run three slightly different copies of the same cleaning do-file, and they get three different answers to the same question. The team's consistency problem is a packaging problem, and this final chapter is how to solve it: turn the scripts you repeat into shared, versioned tools, document them so others can use them, distribute them, and archive the whole project so it outlives the engagement.

We move from a do-file that wants to be a command, through help files and distribution, to a single honest section on where Mata and Python each do what Stata's main language cannot, and end on the replication archive. One claim runs underneath: colleagues treat a tool as a standard only after you package it, document it, distribute it, and preserve it. What the book has demanded of every analysis, that others can rerun it and build on it, now falls on the code itself.

\section{When a do-file wants to be a command}
A block of code you have pasted across five projects is a maintenance liability and a source of the three-answers problem, and the fix is to turn it into a command. The move is to extract the repeated logic into an \texttt{.ado} file, a text file that defines a program and, once it sits on Stata's search path, makes its name work as a command, then replace its hard-coded specifics with arguments and generalize it just enough to serve the cases you actually have.\marginnote{Resist over-generalizing on the first pass. A command that tries to handle cases you do not yet have accretes options nobody uses and bugs nobody catches; the rule of three, refactor into a command once you have pasted the block a third time, keeps the abstraction honest.} The book's own example is \texttt{dsload}\index{dsload@\texttt{dsload}}: the block of project controls and path settings that three colleagues had each modified, unified into one command that takes its paths as arguments so everyone runs the same logic.

\begin{appliedexample}[Turning a one-off script into a shared tool]\label{ae:ch15-sharedtool}
\textbf{Scenario.} You have a do-file block that loads project controls and verifies file paths. Three colleagues use slightly different versions, and path errors follow.

\medskip
\textbf{The packaging path.}
\begin{enumerate}
    \item Extract the logic into \texttt{dsload.ado}, replacing hard-coded paths with command arguments.
    \item Write \texttt{dsload.sthlp} in SMCL to document the syntax and options.
    \item Create the \texttt{dsload.pkg} and \texttt{stata.toc} metadata for distribution.
    \item Push to the team's GitHub repository so colleagues install and update with \texttt{net install}\index{net install@\texttt{net install}}.
\end{enumerate}
One command, one behavior, one answer. The consistency problem was a packaging problem all along.
\end{appliedexample}

Before any of those steps, write the help file's syntax block first, as if the command already existed. Sketch \texttt{dsload using \textit{datafile}, controls(\textit{file}) verify} on paper, read it aloud, and ask whether a colleague could guess what each piece does from the line alone. If the line needs six options to say one thing, or the option names only make sense to the person who wrote them, the design is wrong before any code exists\marginnote{Revising a sentence is cheaper than revising a program, one more reason to draft the help-file syntax line before any code exists.}. The \texttt{syntax} line in the ado then becomes a transcription of a decision already made rather than a decision made while typing.

The skeleton of the command is short, and it is worth seeing once because every shared tool has the same four bones. A program is named and version-locked, a \texttt{syntax}\index{syntax@\texttt{syntax}} line parses what the caller typed into named locals, the body does the work, and \texttt{end} closes it. Everything else is elaboration on these lines:
\par
{\footnotesize
\begin{verbatim}
*! dsload 1.0.0  load project controls + verify paths
program define dsload
    version 17.0
    syntax using/, [Controls(string) Verify]
    // `using' holds the data path; `controls'
    // the settings file; `verify' is a flag.
    confirm file `"`using'"'
    if "`controls'" != "" run `"`controls'"'
    use `"`using'"', clear
    if "`verify'" != "" assert _N > 0
end
\end{verbatim}
}
Read the \texttt{syntax} line as a contract with the caller: \texttt{using/} demands a file path after \texttt{using}, the square brackets mark everything inside them optional, and the capitalized letters in \texttt{Controls} and \texttt{Verify} set the shortest abbreviation a user may type. Whatever the parser accepts lands in local macros of the same names, which the comments trace. Save the file as \texttt{dsload.ado} in your personal ado folder (\texttt{sysdir} prints its location under \texttt{PERSONAL}) and the command exists from your next keystroke.
The \texttt{version 17.0}\index{version@\texttt{version}} line is not decoration. It tells Stata to interpret the code by the rules of that release, so a do-file written today still runs the same way under a future Stata that changed a default. That single line is the replicable principle applied to the tool itself: the command's behavior is pinned to a language version, not to whatever Stata happens to be installed on a colleague's laptop three years from now.\marginnote{This backward-compatibility guarantee is a genuine Stata distinction, not a courtesy: code that opens with \texttt{version 8} still runs unchanged decades later. Few statistical languages promise it, which is one practical reason a reproducibility-minded shop stays in Stata.}

The rule of three says when to generalize; just as useful is a rule for what to leave out. Options nobody asked for are the most common weight a young command carries: an export format added while you were in there, a flag that mirrors a feature admired in someone else's tool, a code path for data shapes the current projects never produce. Each one is syntax to document, a branch to test, and behavior to maintain, so the working rule is that an option enters the command only when some real call in your own projects needs it. And when a second project's quirk would twist the first project's logic, consider forking a sibling command rather than parameterizing: two small commands that each do one thing cleanly outlive one command with a personality split.

Once you package a habit, the team inherits it: colleagues can now extend your work across people, not just across your own projects. A shared command is the difference between everyone doing the same thing and everyone doing almost the same thing. For an evaluator embedded in a nonprofit or agency team, this is not software-engineering for its own sake: it is the same discipline a performance-measurement lead already applies to a metric definition, written once, versioned, and shared so that two analysts computing the same indicator get the same number.\marginnote{\textcite{wilson2017good} make the case that most researchers are never taught these habits and pay for it in lost work and irreproducible results; their remedy is deliberately modest, the \enquote{good enough} practices, versioned code, documented steps, a run order, that a two-person shop can sustain rather than the full apparatus of a software team.}

\section{Help files people actually read}
A tool without documentation dies with its author's memory, so every shared command gets a help file. Stata help files are written in SMCL\index{SMCL}, the Viewer's markup, and one rule keeps them useful: give every command a runnable example, because most people learn a tool by running its example and adapting it.\sidenote{SMCL (\enquote{smickle}, Stata Markup and Control Language) is directive-based: \texttt{\{cmd:...\}} for commands, \texttt{\{help name\}} for cross-links, \texttt{\{synoptset\}} for the options table. It renders only in the Viewer, so a \texttt{.sthlp} file read as plain text looks like tag soup; that is expected.} The \texttt{editanything}\index{editanything@\texttt{editanything}} command, which opens any text file, an ado, a help file, a do-file, from the command line, is the small convenience that lowers the friction of writing and maintaining that documentation.

A help file does not have to be long to be complete. Eight lines carry the parts a reader needs: a title, a syntax line, one sentence of description, and one example they can run. The SMCL below is the whole spine of \texttt{dsload.sthlp}, and it renders in the Viewer as a formatted, cross-linked page (preview a draft any time, no install needed, with \texttt{view dsload.sthlp}):
\par
{\footnotesize
\begin{verbatim}
{smcl}
{title:Title}
{p}{cmd:dsload} -- load project controls, verify
{title:Syntax}
{p}{cmd:dsload using} {it:datafile}{cmd:,}
     [{cmdab:c:ontrols(}{it:file}{cmd:)} {cmd:verify}]
{title:Description}
{p}{cmd:dsload} loads a dataset after running a
 project controls file and confirming its path.
{title:Example}
{p}{cmd:. dsload using data.dta, verify}
\end{verbatim}
}
The example is a command a reader can paste and run, not a description of one. That is what a help file is for: a colleague who runs the example and edits the filename has used the tool without ever opening its source.\marginnote{Run \texttt{help} on any built-in Stata command and copy its shape: title, syntax, description, options, examples, then \enquote{Author.} StataCorp's own help files are the style guide, and matching them makes your tool feel native rather than bolted on.} The syntax block also documents the abbreviation, \texttt{c:ontrols} means \texttt{c} is the shortest legal spelling, so the help file and the \texttt{syntax} line tell the same story.

Written after the code, the help file is documentation; drafted first, it is a specification. Sketching \texttt{dsload.sthlp} before \texttt{dsload.ado} forces the decisions that matter, what the required argument is, which options abbreviate, what the one canonical example looks like, while they still cost nothing to change. A useful check on the draft: hand the syntax line to a colleague and ask what they expect the command to do, because wherever their guess diverges from your intent, the design rather than the reader needs revising.

\begin{kaobox}[frametitle=Package Integration: \texttt{editanything} and \texttt{writeinput}]
\textbf{Integrated commands:} \texttt{editanything} opens any text file from the Stata command line for editing, resolving it across the ado-path and refusing to overwrite base files; \texttt{writeinput}\index{writeinput@\texttt{writeinput}} turns an in-memory dataset into a reproducible \texttt{input} block for tests and bug reports.

The second one solves a problem every bug report has. Asking for help means handing over data, and the data is usually the part you cannot send. \texttt{writeinput} writes the rows themselves as code, so a self-contained example travels in a plain-text message:
\par
{\footnotesize
\begin{verbatim}
. sysuse auto, clear
. keep in 1/2
. keep make price mpg
. writeinput make price mpg using mwe.do, replace

* mwe.do now contains:
clear
input str18 make int price int mpg
`"AMC Concord"' 4099 22
`"AMC Pacer"' 4749 17
end
\end{verbatim}
}
Paste that block into a Statalist post or a test battery and the reader reproduces your exact data before touching your problem. It is also how the test do-files in this chapter get their fixtures: a failing case becomes a permanent test the moment you can write it down.\marginnote{The default output is already round-trip exact: doubles are written with \texttt{\%21.0g} and everything else with \texttt{\%12.0g}, so the pasted block reproduces your values bit for bit. Reach for \texttt{precision(hex)} when you want the bits shown explicitly, as \texttt{\%21x}, so a reader can see at a glance that a \texttt{double} is the issue.}
\end{kaobox}

Hand a colleague a help file with a working example and they can use the tool without ever reading its source; the people who inherit your code get accessibility whether or not anyone names the principle. Beginners skip the documentation step, and it is the step that decides whether a tool survives its author.\marginnote{A house style worth borrowing wholesale: \textcite{cox2002lists} shows, in the \emph{Speaking Stata} column that has taught a generation of user-programmers, how to write loops and list-handling that read cleanly enough for a stranger to follow. A shared command in that idiom feels native rather than improvised.}

\section{Distributing through GitHub and net install}
How does a personal fix become the team standard? Colleagues have to be able to install it, and Stata's package machinery makes that a one-line affair. A \texttt{.pkg} manifest and a \texttt{stata.toc} in a GitHub repository let colleagues and clients install and update a tool with a single \texttt{net install} line they can paste, and versioning the repository means everyone can pin to, or upgrade from, a known release.\marginnote{Adopt semantic versioning for the \texttt{.pkg}: bump the major number only when you break an existing call, the minor for backward-compatible features, the patch for fixes. A user who installed \texttt{net install ..., from(...)} at v1.2 then knows a jump to v2.0 may need their do-files revised, while v1.3 will not.} For machines with no internet, a local SSC catalog (\texttt{scrapessc}\index{scrapessc@\texttt{scrapessc}}) lets an air-gapped analyst still find and install community packages.

Two small text files do the distributing. The \texttt{stata.toc} is the table of contents Stata reads first: it lists which packages a repository offers. The \texttt{.pkg} is the manifest for one package: a title line, a distribution date, and a \texttt{f} line for each file to copy. Together they are what a \texttt{net install} obeys:
\par
{\footnotesize
\begin{verbatim}
* --- stata.toc ---
v 3
d Applied Evaluation tools (Booth & Teas)
p dsload   load controls, verify project paths
* --- dsload.pkg ---
v 3
d dsload: load project controls and verify paths
d Distribution-Date: 20260704
f dsload.ado
f dsload.sthlp
\end{verbatim}
}
The install line a colleague pastes points \texttt{net install} at the raw repository folder, giving the package name and a \texttt{from()} URL. Stata reads \texttt{stata.toc}, finds \texttt{dsload}, reads \texttt{dsload.pkg}, and copies the files into the user's ado-path, the same search path Chapter~\ref{ch:workshop} taught you to stack. Nothing about this needs a package server or an SSC submission; a plain GitHub repository is a working distribution channel, which is why a two-person evaluation shop can ship tools the same way a large group does.\marginnote{SSC, the Boston College archive Baum has curated since 1998, remains the front door for discovery: a package there is one \texttt{ssc install} away and shows up in \texttt{search}. GitHub is faster to publish and iterate; the common path is to develop on GitHub and submit to SSC once the tool has stabilized.}

The version number carries an obligation as well as a signal. Once a colleague's project pins v1.2, that string names a specific behavior, and republishing different behavior under the same number breaks the one guarantee versioning offers, quietly, in their results rather than yours. The discipline is mechanical: any change a user could observe, a new default, a renamed option, a fix that alters output, ships under a new number, and the old release stays retrievable as a tagged commit, a bookmarked point in the repository's history a burned project can always reinstall from. A team that trusts your numbers will pin them; a team burned once will vendor a private copy of the code and never update again, which is the three-answers problem reborn at the package level.

This step widens the tool's reach: a tool nobody can install is a tool nobody uses, so when you publish it you let the rest of the team, and the wider community, benefit from work you did once. The audience you are making the work accessible to has changed, from the reader of a report to a fellow evaluator with an ado-path.

\section{One tool, all the way through the checklist}\label{sec:rateshrink-walk}
The pieces so far, the \texttt{.ado} skeleton, the help file, the two distribution files, are easier to believe as a whole than one at a time, so it helps to watch a single tool cross every line of the checklist. The packages built alongside this book are that worked instance: each began as a one-off block in some engagement and was hardened into a versioned, tested, documented, installable command. Take \texttt{rateshrink}\index{rateshrink@\texttt{rateshrink}}\marginnote{\texttt{rateshrink} is the empirical-Bayes rate stabilizer introduced in Chapter~\ref{ch:trust}; the kaobox later in this section restates what it does.}, and follow it through the same seven steps a colleague would use to judge whether any community tool is safe to depend on.

\begin{figure*}[htbp]
\centering
\begin{tikzpicture}[node distance=5mm and 6mm]
  \node[io] (script) {one-off\\do-file block};
  \node[stage, right=of script] (ado) {\texttt{.ado} +\\\texttt{syntax}};
  \node[stage, right=of ado] (help) {\texttt{.sthlp}\\help file};
  \node[stage, right=of help] (test) {test\\battery};
  \node[stage, below=9mm of test] (ver) {version\\line};
  \node[stage, left=of ver] (pkg) {\texttt{.pkg} +\\\texttt{.toc}};
  \node[io, left=of pkg] (install) {\texttt{net install}\\(shared tool)};

  \draw[flow] (script) -- (ado);
  \draw[flow] (ado) -- (help);
  \draw[flow] (help) -- (test);
  \draw[flow] (test) -- (ver);
  \draw[flow] (ver) -- (pkg);
  \draw[flow] (pkg) -- (install);

  \node[annot, above=0.5mm of ado] {generalize};
  \node[annot, above=0.5mm of help] {document};
  \node[annot, above=0.5mm of test] {prove};
  \node[annot, below=0.5mm of pkg] {distribute};
\end{tikzpicture}
\caption[The hardening path from script to installable tool]{Follow the arrows to see a pasted block become a shared command: it is generalized into an \texttt{.ado} with a \texttt{syntax} line, documented with a help file, proven with a test battery, pinned with a version line, then distributed through the two manifest files so a colleague installs it with one \texttt{net install}. These are the same steps the seven-row checklist audits.}
\label{fig:ch15-viz-harden}
\end{figure*}

The checklist a maintainer owes a user has seven rows, and \texttt{rateshrink} carries all seven. It ships a \emph{help file} (\texttt{rateshrink.sthlp}) with a runnable example, so a reader learns the command by pasting rather than by reading source. It opens with a \emph{version line} in the \texttt{.ado} header, \texttt{*! version 0.1.1  06jul2026}, that pins both the release and the language rules the code was written under. It carries a \emph{test battery} (\texttt{test\_rateshrink.do}) whose six blocks assert that the shrunken rates fall in the unit interval, that reliability rises with the denominator, and that an unsupported option combination exits with an error rather than a wrong answer. It supplies the \emph{\texttt{.pkg} and \texttt{stata.toc}} manifests that let \texttt{net install} copy the right files. It ships a \emph{README} that states what the command does, when to reach for it, and where the method comes from. And it has an \emph{install URL}, a raw GitHub folder a colleague pastes once. The seventh row is the one beginners treat as optional and maintainers treat as load-bearing: a \emph{semantic version} that tells a user what a version jump means before they upgrade.\marginnote{The version string \texttt{0.1.1} is a promise in three numbers \parencite{prestonwerner2013semver}: the leading zero says the public interface may still change, the middle digit counts backward-compatible features, and the trailing digit counts fixes. When \texttt{rateshrink} added the \texttt{rel()} reliability option it moved from a patch to a minor bump, signalling \enquote{new capability, old calls still work} without a word of prose.}

\begin{table*}[htbp]
\centering
\caption[The seven-row tool-hardening checklist]{The seven rows a colleague reads, in order, to decide whether a shared tool is safe to depend on, with the artifact that fills each row and the guarantee it buys. Read it top to bottom against any candidate command: a tool is done when every row is filled.}
\label{tab:ch15-viz-checklist}
\begin{tabular}{lll}
\hline
Row & Artifact & What it guarantees \\
\hline
Help file     & \texttt{rateshrink.sthlp}     & Learn by pasting, not source \\
Version line  & \texttt{*! version 0.1.1}     & Behavior pinned to language \\
Test battery  & \texttt{test\_rateshrink.do}  & Fails loudly if math drifts \\
Manifests     & \texttt{.pkg} + \texttt{.toc} & \texttt{net install} copies files \\
README        & method + limits              & When to use, where it stops \\
Install URL   & raw GitHub folder            & One paste to install \\
Semantic ver. & \texttt{major.minor.patch}    & What a version jump means \\
\hline
\end{tabular}
\end{table*}

\subsection{The two rows that decide whether to depend on a tool}\label{sec:ch15-trust-rows}
The test battery and the citable method are where a shared tool earns trust rather than merely claiming it. The first is the test battery. A command that computes a statistic nobody can check by eye, a shrunken rate is a weighted blend of a raw rate and a prior, needs a test that fails loudly when the arithmetic drifts, because the alternative is a dashboard that shows plausible-but-wrong numbers to a board. The second is the citable method. \texttt{rateshrink}'s help file, README, and \texttt{.ado} header all point at the same companion methods paper, so a reader who wants to know why the reliability weight is defined the way it is can follow one trail from the command to the statistics.\marginnote{Naming the software and the method it implements so both can be cited and found does for one small tool what \textcite{smith2016software} asked of all research code: their six software-citation principles are importance, credit and attribution, unique identification, persistence, accessibility, and specificity.}

The battery is also worth writing before the command it tests. The six blocks in \texttt{test\_rateshrink.do} read as a behavioral contract, shrunken rates stay in the unit interval, reliability rises with the denominator, an unsupported combination exits with an error instead of a number, and each was written down before the arithmetic existed, so the coding session had a finish line: run the battery, watch the failures shrink, stop when the last \texttt{assert} passes. Working in that order turns the question of whether the tool is done from a feeling into a fact, and it leaves the specification behind as a permanent tripwire for every future edit.

The \enquote{weighted blend} the test guards is a single line of arithmetic, and seeing it makes both the test and the reliability column legible. For unit $i$ with raw rate $r_i$ and denominator $n_i$, the shrunken rate is
\[
  \hat{p}_i \;=\; w_i\,r_i \;+\; (1-w_i)\,\bar{p},
  \qquad
  w_i \;=\; \frac{n_i}{n_i + \kappa},
\]
where $\bar p$ is the pooled mean, $w_i$ the per-unit reliability that \texttt{rel()} returns, and $\kappa$ the prior strength the command estimates by method of moments. Section~\ref{sec:shrink} works the intuition behind that weight; what matters here is that $w_i$ \emph{is} the reliability column, which gives the test battery something exact to assert: the weight rises with $n_i$ and stays inside the unit interval.\index{empirical Bayes shrinkage}

\begin{kaobox}[frametitle=Package Integration: \texttt{rateshrink}]
\textbf{Integrated command:} \texttt{rateshrink}\index{rateshrink@\texttt{rateshrink}} applies empirical-Bayes shrinkage to rates computed on small denominators, pulling each unit's raw rate toward the pooled mean in proportion to how little data supports it. It is the worked instance of this chapter's thesis: a one-off shrinkage block, hardened into a versioned, tested, documented, installable tool.

\medskip
\textbf{Install (house idiom).}
\par
{\footnotesize
\begin{verbatim}
local gh "https://raw.githubusercontent.com/ericabooth"
net install rateshrink, ///
    from("`gh'/rateshrink-stata-public/main/") replace force
\end{verbatim}
}

\textbf{Worked passage.} Fifty sites, each with a small and uneven number of observations, report a raw success rate. The command shrinks those rates and returns a per-unit reliability with \texttt{rel()}:
\par
{\footnotesize
\begin{verbatim}
clear
set seed 20260706
set obs 50
gen n = 5 + int(495*runiform()^2)
gen y = rbinomial(n, rbeta(8, 32))
rateshrink, success(y) denominator(n) ///
    generate(pshrunk) ci(95) rel(reliab)
display r(meanrel)
\end{verbatim}
}
On this seeded run the mean reliability across the fifty units is \texttt{r(meanrel)} $= 0.480$, and the largest single move, \texttt{r(maxmove)} $= 0.206$, lands on a site whose raw rate was zero: with only a handful of observations behind it, that zero is pulled up toward the pooled mean of about 0.21 rather than reported as a true zero. The reliability column makes the logic legible, a 208-observation site earns a reliability near 0.65 while a 9-observation site sits near 0.07, so a reader can see at a glance which rates carry weight and which are mostly borrowed from the prior.
\end{kaobox}

Two caveats keep the claim honest and inside what the tool was tested to do. The prior is estimated by method of moments from the same units being shrunk, so with fewer than about ten units the prior itself is noisy and the command falls back to a uniform prior and says so; and the posterior intervals treat that estimated prior as known exactly, so when units are few the intervals run slightly narrower than the true uncertainty. The \texttt{rel()} reliability is defined for the Beta-binomial model only, and asking for it under the Poisson-gamma model exits with an error rather than returning a misleading number.\marginnote{Boundaries like these are what a good README states plainly: a tool that documents where it stops being trustworthy is more useful than one that pretends to have no edges.}

The point of the walk-through is not \texttt{rateshrink} in particular but the pattern it instances. Every shared tool a small evaluation shop ships, whether it stabilizes rates, builds a project skeleton, or suppresses small cells for disclosure review, passes through the same seven rows, and a colleague deciding whether to trust it reads those rows in the same order. When you build your own, the checklist is the specification: help file, version line, test battery, manifest files, README, install URL, semantic version, and the tool is done when every row is filled.

\section{A dash of Mata and Python where it counts}
When should you leave Stata's main language? An evaluator needs the seams between languages, not fluency in all of them, and those seams are few enough to cover honestly in one section. Stata is the system of record and the project manager; Python handles what Stata does not do natively, web scraping, API calls, PDF parsing, LLM requests, as earlier chapters showed; and Mata\index{Mata}, Stata's compiled matrix language, is worth reaching for only where a matrix computation or a tight loop is genuinely too slow in ordinary Stata code.\marginnote{Before rewriting a slow loop in Mata, check whether the slowness is the loop or the disk: a vectorized \texttt{gen}, or Correia's \texttt{gtools} for by-group work, often beats a hand-tuned Mata routine and stays readable. Mata pays off for genuine matrix algebra, not for looping over observations.} The command \texttt{lstrfun}\index{lstrfun@\texttt{lstrfun}}, which manipulates macros with Mata's fast string functions, is the kind of narrow, real win that justifies dropping down a level.

\subsection{A worked example: three ways to build one variable}\label{sec:bench}
The advice to reach for a loop last is easy to give and easy to ignore, so it is worth measuring. The task is deliberately plain: compute a new variable, the squared value of \texttt{x}, on a simulated one-million-row dataset, three ways. The first way is a vectorized \texttt{gen}, Stata's native idiom. The second drops into Mata, reads the column into a matrix, squares it, and writes it back. The third is the anti-pattern: an explicit \texttt{forvalues} loop that touches one observation at a time with \texttt{replace ... in}. We time each with Stata's \texttt{timer}\index{timer@\texttt{timer}}, which is the honest way to settle a speed argument:
\par
{\footnotesize
\begin{verbatim}
set seed 20260704
set obs 1000000
gen double x = runiform()
timer clear
timer on 1
    gen double y1 = x^2                 // vectorized
timer off 1
gen double y2 = .
timer on 2
    mata: st_store(., "y2", st_data(., "x"):^2)
timer off 2
gen double y3 = .
timer on 3
    forvalues i = 1/`=_N' {
        replace y3 = x^2 in `i'         // loop
    }
timer off 3
timer list
\end{verbatim}
}
The Mata line is the one to decode: \texttt{st\_data(., "x")} pulls the variable into Mata as a column, \texttt{:\^{}2} squares it element by element\marginnote{Mata's colon operators work element-wise, which is why \texttt{:\^{}2} squares the column one element at a time.}, and \texttt{st\_store} writes the result back into \texttt{y2}.
The \texttt{timer list} output settles the argument, and the numbers below are the real result of running this do-file (\texttt{ch15\_bench.do}) on simulated data:

\begin{table*}[htbp]
\centering
\caption[Three ways to build one variable, one million rows]{Wall-clock seconds to compute \texttt{x\textasciicircum 2} on one million simulated rows, three ways, timed with Stata's \texttt{timer} in \texttt{ch15\_bench.do}. Vectorized \texttt{gen} wins; Mata beats the explicit loop by orders of magnitude; the row-at-a-time loop is the method to reach for last.}
\label{tab:matabench}
\begin{tabular}{lr}
\hline
Method                       & Seconds \\
\hline
Vectorized \texttt{gen}      & 0.024 \\
Mata matrix column           & 0.030 \\
Explicit \texttt{forvalues}  & 1.555 \\
\hline
\end{tabular}
\end{table*}

Read the table in ratios, not milliseconds. The vectorized \texttt{gen} and the Mata routine both finish in about three hundredths of a second, close enough that the difference between them would not register in a real pipeline. The explicit loop takes more than a second and a half, roughly sixty-five times longer than the \texttt{gen}, to produce the identical column, and the do-file confirms with an \texttt{assert} that all three columns match to twelve digits. The ratio has a mechanical explanation: \texttt{gen} and Mata each hand the whole million-element vector to compiled code once, while the loop pays Stata's per-command interpreter overhead a million separate times. The lesson is the order of reach the margin note above states: vectorize first, drop to Mata only when the operation is genuinely matrix algebra that \texttt{gen} cannot express, and treat an observation-by-observation loop as the method of last resort.\marginnote{The gap looks modest here because \texttt{replace ... in} is one of Stata's cheaper per-observation commands; make the loop body do real work and the same sixty-five-fold penalty becomes the difference between a report that builds in seconds and one that builds over a coffee break.}

The same restraint that keeps unrequested options out of a command keeps extra languages out of a pipeline. When a slow step seems to demand a Mata rewrite, run a timer first: ten lines on simulated data, like the benchmark above, settle in a minute what a hallway debate about speed never settles, and the verdict is often that the rewrite buys nothing a vectorized line does not already deliver.

\begin{kaobox}[frametitle=Package Integration: \texttt{lstrfun}]
\textbf{Integrated command:} \texttt{lstrfun}. Modifies local macros using Mata's compiled string functions, for fast text manipulation inside loops where ordinary macro handling would drag.
\end{kaobox}

An evaluator learns the seams rather than the whole of each language to keep effort in proportion: their time is better spent on the pipeline than on reimplementing it in a faster language that the work does not need. Reach for Mata or Python where it counts, and nowhere else, and the trifecta stays a help rather than a distraction.

\section{The replication archive}\label{sec:archive}
Years after the engagement ends, someone will ask how a number in the final report was computed, and a workflow that cannot be rerun then has not kept its first principle.\marginnote{The replication standard, that a study should ship enough for someone else to reproduce every number, was argued for social science by \textcite{king1995replication} long before it was routine. \textcite{stodden2018reproducibility} later showed the sobering sequel: even journals that require data and code achieve reproducibility far below 100\%, because an archive is only as good as what actually goes in it.} The archive uses two tiers: a canonical copy with a permanent identifier on a repository like openICPSR\index{openICPSR} \parencite{openicpsr} (a DOI, versioning, longevity), paired with a convenience mirror on GitHub, from which a single \texttt{use} of the file's raw URL loads the data in one line with no login.\marginnote{The division of labour matters: cite the openICPSR DOI in the paper, because a DOI is promised to resolve for decades and freezes an exact version, whereas a GitHub raw URL is a convenience that can move, rename, or vanish the day someone force-pushes. Never cite the raw URL as the archival record.} The canonical copy is for permanence and citation; the mirror is for the colleague who just wants to run the code today.

What goes into the archive is not a folder dump but a manifest, so a stranger can tell at a glance that nothing is missing. The manifest below is the one this book's own replication archive\index{replication archive} uses, and it doubles as a checklist: every row names an item, points at where the book demonstrated it, and says why a replicator needs it.

\begin{table*}[htbp]
\centering
\caption[A replication-archive manifest]{The replication-archive manifest for this book's own materials. Each row is an item a stranger needs to reproduce the work without asking, the chapter or command that demonstrates it, and the reason it belongs in the archive.}
\label{tab:manifest}
\begin{tabular}{lll}
\hline
Item              & Example              & Why it belongs \\
\hline
Raw data          & QCEW CSVs (Ch.~\ref{ch:apis})  & The unedited source \\
Cleaning code     & \texttt{ch03\_*.do}  & Raw to analytic file \\
Analysis code     & \texttt{ch07\_trust.do} & Tables and figures \\
Codebook          & \texttt{writeinput}  & What each variable is \\
Shared tools      & \texttt{dsload.ado}  & Custom commands used \\
Environment note  & \texttt{version 17.0} & Stata + package state \\
Read-me           & \texttt{00\_control.do} & Run order, one entry \\
DOI record        & openICPSR            & Permanent citation \\
\hline
\end{tabular}
\end{table*}

The manifest is worth assembling because each row answers a question a replicator would otherwise have to email you. The environment note is the row beginners skip and regret: without a record of which Stata version and which community packages produced the tables, a rerun three years later can silently differ, and \texttt{version 17.0} in the code plus a one-line list of installed packages is enough to close that gap.\marginnote{Pin shared tools too: record the exact release of every house command the code calls, \texttt{rateshrink 0.1.1} rather than \texttt{rateshrink}, because a replicator who installs today's release may be running behavior the original analysis never saw.} Into the archive go the data, the code, the codebooks, and a record of the environment, everything a stranger would need to reproduce the work without asking. That is the four principles kept after the contract ends: the analysis stays replicable, the next team can extend it, a reader can access it, and the findings remain something someone can act on. An evaluation that is archived this way is one whose credibility does not expire with its funding.

\paragraph{Where this leaves you.} By now, you can turn a repeated do-file into a documented, distributed command, complete with a version-locked \texttt{.ado} skeleton, an eight-line help file, and the two-file \texttt{.pkg} plus \texttt{stata.toc} that make it a one-line install. In real timer output, you have seen why you reach for a vectorized \texttt{gen} first, Mata only for genuine matrix work, and an explicit loop last. Beyond the command itself, you can package a project into a manifest-driven replication archive that survives the engagement. Those are extensibility and replicability applied to the practice itself, so the work outlasts any single project. And \texttt{rateshrink}, one of the book's own packages, served as the worked case: you walked it through the seven-row checklist a colleague uses to decide whether to depend on a shared tool, the template for hardening any of your own scripts the same way. The habits underneath the checklist travel too: draft the syntax line and the test battery before the code, generalize only what a real call site needs, and never change behavior under a version number a colleague has pinned. The appendices that follow collect the setup guide, the causal quick-reference, the full toolkit, two complete worked examples, a capstone you can run end to end in minutes, and a registry of every dataset the book uses.

\section*{Exercises}\addcontentsline{toc}{section}{Exercises}
\begin{enumerate}
    \item \emph{Try it on your own data.} Find a code block you have pasted into three or more of your own do-files, extract it into an \texttt{.ado} file that opens with \texttt{version 17.0} and takes its paths as arguments, and confirm the tool works by running your own analysis through the new command and checking that the numbers match the pasted version.
    \item Write an eight-line \texttt{dsload.sthlp} help file for the \texttt{dsload} skeleton in this chapter, giving it a title, a syntax line, a one-sentence description, and one runnable example, then open it with \texttt{help dsload} and confirm the Viewer renders the example as a command you can paste.
    \item Build the two-file distribution for \texttt{dsload}, a \texttt{stata.toc} and a \texttt{dsload.pkg} that lists the \texttt{.ado} and \texttt{.sthlp} files, push them to a GitHub repository, and verify the package installs by running \texttt{net install} from the raw repository URL on a second machine or a clean ado-path.
    \item Rerun \texttt{ch15\_bench.do} after raising the row count to five million, predict how many times slower the explicit \texttt{forvalues} loop will run than the vectorized \texttt{gen} before you read \texttt{timer list}, then compare your prediction against the ratio the timer reports.
    \item Break the benchmark on purpose: change the loop body in \texttt{ch15\_bench.do} so it computes \texttt{x\textasciicircum 3} instead of \texttt{x\textasciicircum 2}, rerun the do-file, and confirm that the \texttt{assert} comparing the three columns stops the run rather than letting the mismatched result pass silently.
    \item \emph{Read a real tool against the checklist.} Install \texttt{rateshrink} with the house install idiom, then open its help file, README, and \texttt{.ado} header and confirm each of the seven checklist rows, help file with a runnable example, version line, test battery, \texttt{.pkg}/\texttt{stata.toc}, README, install URL, and a semantic version, is present; then run the worked example from this chapter and check that \texttt{r(meanrel)} and \texttt{r(maxmove)} match the values reported here.
    \item \emph{Design before you build.} Pick a block you intend to turn into a command, draft its help-file syntax block and a five-assert test battery before touching the \texttt{.ado}, then code until the battery passes, noting which design decisions the syntax draft forced you to settle early.
\end{enumerate}
